% ======================================================================
% DOCUMENTCLASS
% ======================================================================
\documentclass[12pt,oneside]{book}

% ======================================================================
% METADATOS
% ======================================================================
\newcommand{\universidad}{Universidad de Buenos Aires}
\newcommand{\materia}{Derecho Civil}
\newcommand{\titulo}{Ausencia y Presunción de Fallecimiento}
\newcommand{\autor}{Isaac Edgar Camacho Ocampo}
\newcommand{\anio}{2025}

% ======================================================================
% PAQUETES
% ======================================================================
\usepackage[utf8]{inputenc}
\usepackage[T1]{fontenc}
\usepackage[spanish,es-nodecimaldot]{babel}

\usepackage[
    top=2.5cm,
    bottom=2.5cm,
    left=3cm,
    right=2.5cm,
    headheight=15pt
]{geometry}

\usepackage{amsmath}
\usepackage{amssymb}
\usepackage{mathtools}

\usepackage{graphicx}
\usepackage{float}
\usepackage{caption}
\usepackage{subcaption}

\usepackage{booktabs}
\usepackage{array}
\usepackage{longtable}
\usepackage{tabularx}

\usepackage{enumitem}

\usepackage{xcolor}
\usepackage[most]{tcolorbox}

\usepackage{fancyhdr}

\usepackage{ragged2e}

\usepackage[
    colorlinks=true,
    linkcolor=blue!50!black,
    citecolor=green!40!black,
    urlcolor=blue!60!black,
    pdftitle={\titulo},
    pdfauthor={\autor},
    pdfsubject={Apuntes universitarios},
    bookmarks=true
]{hyperref}

% ======================================================================
% RUTAS DE IMÁGENES
% ======================================================================
\graphicspath{
    {./img/}
    {./graficos/}
    {./figuras/}
}

% ======================================================================
% CONFIGURACIÓN
% ======================================================================
\setcounter{tocdepth}{2}

\setlength{\parindent}{0pt}
\setlength{\parskip}{0.5em}

\setlist[itemize]{
    topsep=0.3em,
    itemsep=0.2em
}

\setlist[enumerate]{
    topsep=0.3em,
    itemsep=0.2em
}

\clubpenalty=10000
\widowpenalty=10000

% ======================================================================
% COMANDOS
% ======================================================================
\newcommand{\abs}[1]{\left|#1\right|}
\newcommand{\norm}[1]{\left\lVert#1\right\rVert}
\newcommand{\vect}[1]{\mathbf{#1}}
\newcommand{\deriv}[2]{\frac{d#1}{d#2}}
\newcommand{\pderiv}[2]{\frac{\partial#1}{\partial#2}}

% ======================================================================
% ENTORNOS / CAJAS
% ======================================================================
\newtcolorbox{definition}[1]{
    enhanced,
    breakable,
    title={Definición: #1},
    fonttitle=\bfseries,
    colback=blue!3,
    colframe=blue!50!black,
    boxrule=0.8pt,
    arc=2mm
}

\newtcolorbox{important}{
    enhanced,
    breakable,
    title={Importante},
    fonttitle=\bfseries,
    colback=yellow!8,
    colframe=orange!70!black,
    boxrule=0.8pt,
    arc=2mm
}

\newtcolorbox{example}[1]{
    enhanced,
    breakable,
    title={Ejemplo: #1},
    fonttitle=\bfseries,
    colback=green!4,
    colframe=green!40!black,
    boxrule=0.8pt,
    arc=2mm
}

\newtcolorbox{formula}[1]{
    enhanced,
    breakable,
    title={Fórmula / Regla: #1},
    fonttitle=\bfseries,
    colback=purple!4,
    colframe=purple!50!black,
    boxrule=0.8pt,
    arc=2mm
}

\newtcolorbox{exercise}[1]{
    enhanced,
    breakable,
    title={Ejercicio: #1},
    fonttitle=\bfseries,
    colback=gray!5,
    colframe=gray!60!black,
    boxrule=0.8pt,
    arc=2mm
}

\newtcolorbox{note}{
    enhanced,
    breakable,
    title={Nota},
    fonttitle=\bfseries,
    colback=white,
    colframe=gray!50,
    boxrule=0.6pt,
    arc=2mm
}

% ======================================================================
% CONFIGURACIÓN FINAL (encabezados y pies)
% ======================================================================
\pagestyle{fancy}
\fancyhf{}

\fancyhead[L]{\nouppercase{\leftmark}}
\fancyhead[R]{\materia}

\fancyfoot[C]{\thepage}

\renewcommand{\headrulewidth}{0.4pt}
\renewcommand{\footrulewidth}{0pt}

\fancypagestyle{plain}{
    \fancyhf{}
    \fancyfoot[C]{\thepage}
    \renewcommand{\headrulewidth}{0pt}
}

% ======================================================================
% DOCUMENT
% ======================================================================
\begin{document}

% ----------------------------------------------------------------------
% PORTADA
% ----------------------------------------------------------------------
\begin{titlepage}

\thispagestyle{empty}

\begin{center}

\vspace*{1cm}

{\large \textsc{\universidad}}

\vspace{3cm}

{\Huge \textbf{\materia}}

\vspace{1cm}

{\LARGE \titulo}

\vspace{0.8cm}

\textit{Comprender los conceptos es el primer paso para convertir el estudio en conocimiento.}

\vspace{1.2cm}

\rule{0.70\textwidth}{0.5pt}

\vspace{0.5cm}

{\large Apuntes de estudio}

\vspace{0.5cm}

\rule{0.70\textwidth}{0.5pt}

\vfill

{\large \autor}

\vspace{0.3cm}

{\large \anio}

\end{center}

\vfill

\begin{flushleft}
\textit{Este es un modesto aporte a los estudiantes de ``Derecho Civil'' no pretende sustituir el material oficial de la materia.}
\end{flushleft}

\end{titlepage}

% ----------------------------------------------------------------------
% ÍNDICE
% ----------------------------------------------------------------------
\tableofcontents

% ----------------------------------------------------------------------
% PRESENTACIÓN DEL TEMA
% ----------------------------------------------------------------------
\chapter*{Presentación del tema}
\addcontentsline{toc}{chapter}{Presentación del tema}

Estos apuntes abordan el fin de la existencia de la persona en su vertiente
de incertidumbre sobre su vida o muerte, regulado en el Código Civil y
Comercial de la Nación (CCCN) a través de dos institutos diferenciados: la
\textbf{ausencia} (Arts. 79-84 CCCN) y la \textbf{presunción de
fallecimiento} (Arts. 85-92 CCCN).

Los temas tratados son los siguientes:

\begin{itemize}
    \item Fin de la existencia de la persona: marco conceptual y legal.
    \item Distinción entre ausencia (Arts. 79-84 CCCN) y presunción de fallecimiento (Arts. 85-92 CCCN).
    \item Antecedentes históricos: Vélez Sársfield, Ley 14.394 y el nuevo CCCN.
    \item Casos ordinario, extraordinario genérico y extraordinario específico.
    \item Procedimiento judicial, sentencia y día presuntivo de fallecimiento.
    \item Garantías para el caso de reaparición y prenotación registral.
\end{itemize}

\begin{note}
Como abogado o abogada, tarde o temprano resulta habitual encontrarse con un
cliente cuyo familiar desapareció: un accidente en ruta, una catástrofe
natural, o una persona que simplemente dejó de dar señales de vida.
\textbf{Saber distinguir la ausencia de la presunción de fallecimiento no es
un detalle académico: es la diferencia entre proteger el patrimonio familiar
desde el primer mes o perderlo por falta de acción.}

Estos institutos permiten: (1) asesorar a herederos sobre cuándo y cómo
iniciar acciones judiciales; (2) gestionar la curatela de bienes de personas
desaparecidas; (3) representar acreedores que necesitan cobrar deudas
pendientes; (4) intervenir en procesos sucesorios abiertos tras una
declaración de fallecimiento presunto; (5) proteger los derechos del
ausente ante una eventual reaparición, gestionando la prenotación
registral.

En síntesis, estas reglas existen para que el derecho privado pueda dar una
respuesta ordenada, justa y temporalmente acotada a uno de los dramas más
dolorosos de la vida real: no saber si un ser querido está vivo o muerto.
\end{note}

\textbf{¿Cómo se relacionan los temas?} Puede pensarse el proceso como una
escalada progresiva: si no se tienen noticias de una persona y sus bienes
quedan desprotegidos, corresponde solicitar primero la \textbf{declaración
de ausencia} (Arts. 79-84) para que un curador los cuide. Si transcurren
tres años sin novedades, puede iniciarse la \textbf{presunción de
fallecimiento} (Arts. 85-92), que abre la sucesión pero con garantías por
si la persona reaparece. Todo el proceso protege tanto a quienes esperan
como a quien podría volver.

% ======================================================================
% CAPÍTULO 1
% ======================================================================
\chapter{Marco conceptual e histórico}

\section{Introducción: ausencia y presunción de fallecimiento}

El fin de la existencia de la persona está regulado en el Código Civil y
Comercial a través de dos figuras: la \textbf{ausencia}, por un lado, y la
\textbf{presunción de fallecimiento}, por el otro.

Ausencia y presunción de fallecimiento son dos figuras distintas. Es
habitual que se confundan bajo la expresión «ausencia con presunción de
fallecimiento», lo cual no resulta exacto: se trata de institutos
diferentes. Por supuesto, la presunción de fallecimiento supone que existe
una ausencia de la persona, pero lo que corresponde estudiar es,
específicamente, por un lado la \textbf{ausencia} y por el otro la
\textbf{presunción de fallecimiento}.

\section{El problema jurídico de fondo}

El problema jurídico de fondo puede formularse como la pregunta: ¿cuál es
la respuesta jurídica ante la ausencia prolongada de una persona de su
domicilio?

Esta cuestión ya encontraba explicación en el contexto de las guerras: la
persona que quedaba cautiva y de la que no se sabía nada. ¿Qué sucedía con
sus bienes? ¿Quién podía tomar medidas conservatorias? También pueden
pensarse situaciones catastróficas más recientes: la desaparición de un
avión, el tsunami que hubo en Tailandia a comienzos del año 2002, un
incendio, o cualquier catástrofe en la que una persona desaparece y no se
sabe si está viva o no.

Se produce una \textbf{incertidumbre}, y esa incertidumbre no puede durar
eternamente, porque para el derecho civil existe una cantidad de
relaciones y situaciones jurídicas que ameritan respuesta y toma de
decisiones.

No existe una institución estatal que actúe de oficio ante la mera falta
de contacto de una persona: este es un \textbf{sistema de derecho civil,
de derecho privado}, en el cual los particulares que tienen algún tipo de
interés en los bienes o en las relaciones jurídicas acuden a un proceso
judicial, ya sea para declarar la ausencia o para declarar la presunción
de fallecimiento.

Esto supone preguntarse: ¿cuándo se considera que una persona ha
fallecido? ¿Cuándo resulta razonable asumirlo así? ¿Quién puede pedirlo?
¿Qué sucede si la persona reaparece? ¿Qué resguardos jurídicos se toman en
el camino por si reaparece? Ese es el tema de fondo que estructura ambos
institutos.

\begin{important}
No debe confundirse la ausencia con la presunción de fallecimiento. Si
bien están ubicadas inmediatamente una al lado de la otra en el Código
Civil y Comercial, y tienen elementos comunes, son institutos distintos.
\end{important}

\section{Antecedentes históricos}

Los antecedentes de esta institución provienen del \textbf{Código de
Vélez Sársfield}, que regulaba la ausencia con presunción de fallecimiento
en los artículos 110 a 625. La doctrina señala con claridad cómo en Vélez
Sársfield se produce una \textbf{confusión de elementos}:

Por un lado, existen antecedentes del \textbf{derecho germánico}, donde
existía claramente el instituto de la declaración de muerte, que Vélez
recibe a través de Freitas. Por el otro lado, el \textbf{derecho
francés}, que no conocía esa declaración de muerte y simplemente regulaba
la ausencia, sin poner fin a la existencia de la persona. Vélez también
reconoce esta figura de la ausencia, de modo que el instituto queda
regulado con elementos mixtos.

Tal como fue regulado por Vélez, la institución finalmente recibió
críticas. En 1984 una ley ómnibus —una ley que reunía muchos temas en su
contenido—, la \textbf{Ley 14.394}, reguló por un lado la ausencia y por
otro lado la presunción de fallecimiento: dos instituciones distintas,
clarificando y resolviendo lo que Vélez había planteado. El régimen de
Vélez, entonces, quedó reemplazado por el sistema de la Ley 14.394.

\begin{important}
\textbf{CCCN — Estructura normativa actual.} El nuevo Código Civil y
Comercial deroga la Ley 14.394 en estas partes, pero sigue sus
lineamientos e incorpora tanto la ausencia como la presunción de
fallecimiento al articulado del Título I, Libro I:
\begin{itemize}
    \item \textbf{Ausencia} — Capítulo 6, Arts. 79-84 (no implica fin de la existencia de la persona).
    \item \textbf{Presunción de Fallecimiento} — Capítulo 7, Arts. 85-92 (implica inscripción en el Registro Civil y apertura de la sucesión).
\end{itemize}
Metodológicamente, el CCCN los incorpora con muy pocos cambios respecto de
la ley anterior.
\end{important}

% ======================================================================
% CAPÍTULO 2
% ======================================================================
\chapter{Ausencia (Arts. 79-84 CCCN)}

\section{Concepto de ausencia}

La ausencia es una institución relativamente sencilla. Si la persona
desapareció de su domicilio sin tenerse noticias de ella y sin dejar
apoderado, o si el apoderado no tiene atribuciones suficientes, puede
asignarse un \textbf{curador a los bienes}.

\begin{definition}{Ausencia}
La ausencia, en el proceso del Art. 79 CCCN, no busca declarar a la
persona presuntamente fallecida. Simplemente busca \textbf{garantizar la
adopción de medidas sobre los bienes} de una persona de la que no se
tienen noticias.
\end{definition}

\begin{example}{¿Qué NO es ausencia?}
Si alguien quedó varado en Qatar durante la pandemia, se tienen noticias
de esa persona: se sabe que está varada y que no puede regresar a su
domicilio. Eso no configura una ausencia. La ausencia es la situación de
la persona que desaparece, no se conecta a internet, no usa su tarjeta de
crédito, no pasa por migraciones, no tiene registro en su correo
electrónico, no tiene ningún tipo de actividad, y por lo tanto nadie de
las personas conocidas tiene noticias de ella.
\end{example}

\section{Legitimación y juez competente}

\begin{formula}{Art. 80 CCCN — Legitimación activa}
El Art. 80 establece quiénes pueden iniciar la acción: el Ministerio
Público y las personas que tengan interés legítimo en los bienes (por
ejemplo, los acreedores o los sucesores).
\end{formula}

Respecto del juez competente:

\begin{table}[H]
\centering
\begin{tabularx}{\textwidth}{>{\raggedright\arraybackslash}X >{\raggedright\arraybackslash}X}
\toprule
\textbf{Situación} & \textbf{Juez competente} \\
\midrule
Persona con domicilio en Argentina & Juez del domicilio \\
Sin domicilio en el país o lugar desconocido & Juez del lugar donde están los bienes que hay que cuidar \\
Bienes en distintas jurisdicciones & El primero que intervenga (principio de prevención) \\
\bottomrule
\end{tabularx}
\caption{Juez competente en el proceso de ausencia}
\end{table}

\section{Procedimiento y curatela}

Se cita al ausente por \textbf{cinco días por edictos}. Los edictos son
los avisos que se publican en el Boletín Oficial y en diarios de
circulación importante, mediante los cuales se convoca a la persona a que
se haga presente. Si no comparece, se le da intervención al
\textbf{defensor oficial}, que interviene en representación del ausente,
y pueden tomarse medidas precautorias para los bienes.

La prueba se orienta a demostrar que no hay noticias de esta persona: se
libran distintos oficios, se solicitan testigos, y se determina si existe
o no ausencia. El Ministerio Público procura defender al presunto
ausente, evaluando la verosimilitud de la situación alegada.

Determinada la ausencia, se nombra un curador. El procedimiento del
curador comprende actos de \textbf{conservación y administración de los
bienes}. No puede realizar actos extraordinarios: no puede vender ni
hipotecar un inmueble, porque en ese caso estaría comprometiendo el
patrimonio. Los actos extraordinarios que resulten de necesidad evidente e
impostergable requieren \textbf{autorización judicial}.

Los frutos —por ejemplo, la renta de un inmueble que se alquila— se
destinan al \textbf{sostenimiento de los familiares del ausente}.

\section{Conclusión de la curatela}

La curatela concluye:

\begin{enumerate}[label=(\alph*)]
    \item Porque reaparece el ausente, personalmente o por apoderado.
    \item Porque se constata su muerte: ya sea porque aparece el cadáver y se hace la inscripción del fallecimiento y la correspondiente partida, o bien porque se declara el fallecimiento presunto judicialmente, instituto que se analiza a continuación.
\end{enumerate}

% ======================================================================
% CAPÍTULO 3
% ======================================================================
\chapter{Presunción de fallecimiento (Arts. 85-92 CCCN)}

El sistema del Código es igual al sistema de la Ley 14.394 y se estructura
en torno a \textbf{tres casos}: un caso ordinario y dos casos
extraordinarios —el extraordinario genérico y el extraordinario
específico—.

\section{Caso ordinario (Art. 85 CCCN)}

\begin{formula}{Art. 85 CCCN — Caso ordinario}
La persona que se ausenta de su domicilio sin que se tengan noticias de
ella por el término de \textbf{tres (3) años}. Los elementos son: la
ausencia sin noticias durante tres años. No hay ningún elemento
particular adicional.
\end{formula}

El plazo es de \textbf{tres años}. Esto se vincula con situaciones como,
por ejemplo, una persona que desaparece —quizás víctima de un homicidio
no resuelto— y de la que no se sabe bien qué pasó. Al cabo de tres años
de incertidumbre, los familiares pueden solicitar su fallecimiento
presunto.

No es obligatorio: el cumplimiento del plazo de tres años no impone
iniciar el proceso de manera obligatoria. Esto se relaciona con la
certeza de las relaciones jurídicas. Dado que en la mayoría de los casos
se trata de un familiar y todo está rodeado de afectos y conmociones, hay
quienes siguen buscando y pueden dejar pasar cinco o seis años. Lo que
debe demostrarse en el proceso es que los supuestos legales se han
cumplido.

\begin{important}
\textbf{La última noticia.} El plazo de tres años —o el que corresponda
según el caso— se cuenta desde la \textbf{última noticia} que se tuvo de
la persona. Por ejemplo: si alguien cruzó por migraciones en una fecha
determinada por el paso fronterizo de Paso de los Libres, esa es la
última noticia que se tuvo. A partir de allí se cuentan los plazos.
\end{important}

\section{Caso extraordinario genérico (Art. 86, 1.\textsuperscript{er} párr.)}

El caso extraordinario se denomina así porque acontece algún hecho que
marca una \textbf{mayor probabilidad} de que haya habido un deceso de la
persona. El caso genérico comprende situaciones de incendio, terremoto,
acción de guerra u \textit{otro suceso semejante}, o bien la
participación de la persona en una actividad que implique el mismo
riesgo.

\begin{formula}{Art. 86, 1.\textsuperscript{er} párr. — Caso extraordinario genérico}
Si la persona hubiese estado en el lugar de un incendio, terremoto,
acción de guerra u otro suceso semejante susceptible de ocasionar la
muerte, o hubiese participado de una actividad que implique el mismo
riesgo, y no se tuviese noticias de ella durante \textbf{dos (2) años}
desde el día en que el suceso ocurrió o pudo haber ocurrido.
\end{formula}

\begin{example}{Casos de extraordinario genérico}
\begin{itemize}
    \item El alpinista que estaba escalando el Aconcagua o el Everest y dejó de dar señales.
    \item El tsunami de Tailandia.
    \item Un incendio, un terremoto, una acción bélica.
    \item Cualquier suceso que resulte muy probable que haya provocado la muerte.
\end{itemize}
\end{example}

En este caso, el plazo se adelanta a \textbf{dos años}, contados desde el
día en que el suceso ocurrió o pudo haber ocurrido. Por ejemplo, si se
trata de un alpinista y el último día que consta en un registro que
estaba subiendo fue el 10 de mayo de 2020, desde allí se cuentan los dos
años.

\section{Caso extraordinario específico (Art. 86, 2.\textsuperscript{do} párr.)}

El caso extraordinario específico es el supuesto más detallado por el
Código: el de un \textbf{buque naufragado o aeronave perdida}. Aquí el
acontecimiento resulta mucho más probable como causa de muerte, y el
plazo es de \textbf{seis meses} desde que el suceso ocurrió o pudo
ocurrir.

\begin{formula}{Art. 86, 2.\textsuperscript{do} párr. — Caso extraordinario específico}
Si el presunto fallecido estuvo en un buque naufragado o aeronave
siniestrada o perdida, y no se tuviese noticias de su suerte durante
\textbf{seis (6) meses} desde el día en que el suceso ocurrió o pudo
haber ocurrido.
\end{formula}

\section{Tabla comparativa de los tres casos}

\begin{table}[H]
\centering
\begin{tabularx}{\textwidth}{>{\raggedright\arraybackslash}X c c >{\raggedright\arraybackslash}X}
\toprule
\textbf{Tipo de caso} & \textbf{Artículo} & \textbf{Plazo} & \textbf{Circunstancia} \\
\midrule
Ordinario & Art. 85 & 3 años & Sin noticias \\
Extraordinario genérico & Art. 86, 1.\textsuperscript{er} párr. & 2 años & Catástrofe / actividad de riesgo \\
Extraordinario específico & Art. 86, 2.\textsuperscript{do} párr. & 6 meses & Buque naufragado / aeronave perdida \\
\bottomrule
\end{tabularx}
\caption{Comparación de los tres casos de presunción de fallecimiento}
\end{table}

\section{Casos reales mencionados en clase}

\begin{example}{Caso ``Los Andes'' (1972) — Aeronave perdida}
En 1972, un avión militar uruguayo iba rumbo a Chile y cayó en la mitad
de la cordillera de los Andes. Durante 72 días se buscó a sus ocupantes
hasta que dos supervivientes lograron cruzar hacia el lado chileno y
dieron aviso. Esas personas estuvieron desaparecidas durante 72 días. Si
hubiera existido necesidad de declarar su fallecimiento, habría
correspondido el caso extraordinario específico (buque/aeronave), por lo
que hubiera sido necesario esperar seis meses.
% TODO: contenido incompleto o ambiguo en las notas originales — el apunte indica que "los encontraron antes de los 62 días", cifra que no coincide con los "72 días" mencionados previamente en el mismo párrafo; se conserva tal como figura en el material original.
Como los encontraron antes de los 62 días, ese plazo no llegó a cumplirse.
\end{example}

\begin{example}{Caso Malaysia Airlines MH370 (2014) — Aeronave perdida}
En 2014 desapareció un avión de Malaysia Airlines y nunca se supo con
certeza dónde terminó. El último contacto se toma como fecha para contar
los seis meses. Quien tuviera un familiar en ese vuelo debía esperar, como
mínimo, seis meses desde ese último contacto para poder iniciar el
proceso de presunción de fallecimiento bajo el caso extraordinario
específico.
\end{example}

% ======================================================================
% CAPÍTULO 4
% ======================================================================
\chapter{Procedimiento y sentencia}

\section{Legitimación y juez competente}

\begin{formula}{Art. 87 CCCN — Legitimación activa}
Cualquier persona que tenga un derecho subordinado a la muerte puede
pedir la declaración del fallecimiento presunto y debe justificar los
extremos legales: la ausencia, la falta de noticias y las diligencias
realizadas para averiguar la existencia del ausente.
\end{formula}

El \textbf{juez competente} es el juez del domicilio del ausente. Si bien
las normas procesales corresponden a cada provincia, en este caso el
Código Civil y Comercial las incorpora para dar una
\textbf{regulación común a todo el país}, evitando diferencias en el
tratamiento de una situación muy delicada, ya que se está declarando
fallecida a una persona para iniciar su sucesión.

\section{Diligencias probatorias}

Para justificar los extremos legales deben realizarse diligencias
tendientes a averiguar la existencia del ausente, entre ellas:

\begin{itemize}
    \item Oficios a distintas entidades.
    \item Pedido de informes a su tarjeta de crédito.
    \item Pedido de informes a su cuenta bancaria.
    \item Pedido a la empresa de telefonía celular.
    \item Pedido al club si es socio.
    \item Pedido a Migraciones para verificar entradas o salidas del país.
    \item Pedido de informe sobre actividad de su celular.
\end{itemize}

Se le nombra un \textbf{defensor al ausente} (o el defensor oficial). Se
lo cita por \textbf{edictos una vez por mes durante seis meses}: se
publican seis edictos, uno por mes. Se le nombra, asimismo, un curador de
los bienes si no hay quien los esté administrando.

\begin{important}
\textbf{Aclaración del Art. 79 (ausencia simple) frente a la presunción
de fallecimiento.} El Código aclara que el hecho de haber tramitado
previamente la simple ausencia (Art. 79) \textbf{no} es presupuesto
necesario ni suficiente para la presunción de fallecimiento. Es decir:
(a) puede iniciarse directamente la presunción de fallecimiento sin haber
tramitado la ausencia previa; y (b) haber tramitado la ausencia no exime
de tener que probar nuevamente, en el juicio de presunción de
fallecimiento, que se han realizado todas las diligencias para dar con el
ausente. Ambos procesos son independientes, aunque puedan darse en forma
sucesiva.
\end{important}

\section{Sentencia y día presuntivo de fallecimiento}

Pasados los seis meses de citación, recibida la prueba y oído el
defensor, el juez debe dictar sentencia. La sentencia contiene tres
elementos:

\begin{enumerate}
    \item Declara el fallecimiento presunto, comprobados que estén los extremos legales.
    \item Fija el día presuntivo de fallecimiento.
    \item Dispone la inscripción de la sentencia en el Registro Civil.
\end{enumerate}

\textbf{¿Cómo se fija el día presuntivo?}

\begin{table}[H]
\centering
\begin{tabularx}{\textwidth}{>{\raggedright\arraybackslash}X >{\raggedright\arraybackslash}X >{\raggedright\arraybackslash}X}
\toprule
\textbf{Caso} & \textbf{Día presuntivo} & \textbf{Ejemplo} \\
\midrule
Ordinario (Art. 85) & Último día del 1.\textsuperscript{er} año = mitad del plazo de 3 años & Última noticia: 31/12/2017 $\rightarrow$ día presuntivo: 30/06/2019 \\
Extraordinario genérico (Art. 86, 1.\textsuperscript{er} párr.) & El día en que ocurrió el suceso (o el término medio si se extendió en el tiempo) & Si el incendio duró varios días: el día del mediodía del período \\
Extraordinario específico (Art. 86, 2.\textsuperscript{do} párr.) & El último día en que se tuvo noticia del buque o aeronave & MH370: último contacto = 8 de marzo de 2014 $\rightarrow$ ese es el día presuntivo \\
\bottomrule
\end{tabularx}
\caption{Fijación del día presuntivo de fallecimiento según el caso}
\end{table}

El día presuntivo es de gran relevancia porque determina
\textbf{quiénes son los herederos} que entran en la sucesión del
presuntamente fallecido, declarado judicialmente. Si es posible, la
sentencia debe determinar en qué hora presuntiva se fijó el
fallecimiento; si no, se toma al final de ese día.

% ======================================================================
% CAPÍTULO 5
% ======================================================================
\chapter{Garantías y efectos de la reaparición}

\section{Las tres garantías adoptadas}

Con la declaración de fallecimiento presunto se abre la sucesión y los
herederos pueden tomar los bienes. Pero el Código adopta \textbf{tres
garantías} por si el ausente reapareciera:

\begin{table}[H]
\centering
\begin{tabularx}{\textwidth}{>{\raggedright\arraybackslash}p{0.28\textwidth} >{\raggedright\arraybackslash}X}
\toprule
\textbf{Garantía} & \textbf{Descripción} \\
\midrule
\textbf{Inventario} & Se hace un inventario detallado de todos los bienes del presuntamente fallecido (el acervo sucesorio / masa hereditaria). En base a ese inventario, si reaparece, los herederos deben rendir cuentas de todos esos bienes. \\
\textbf{Prenotación registral} & Al transferir los bienes a nombre de los herederos (por ejemplo, un departamento a nombre de dos hijos), se hace una prenotación: se deja constancia en el asiento registral de que el bien está sometido al régimen de declaración de fallecimiento presunto. \\
\textbf{Restricción para enajenar/gravar} & Para enajenar (vender) o gravar (hipotecar) los bienes heredados, los herederos necesitan autorización judicial. La prenotación avisa a terceros contratantes de que ese bien está sometido a este régimen. \\
\bottomrule
\end{tabularx}
\caption{Las tres garantías del sistema ante una eventual reaparición}
\end{table}

\section{Duración de la prenotación (Art. 91 CCCN)}

\begin{formula}{Art. 91 CCCN — Prenotación}
La prenotación caduca a los \textbf{cinco (5) años} desde el día
presuntivo del fallecimiento, o cuando la persona hubiera cumplido
\textbf{ochenta (80) años} de edad, lo que ocurra primero. Vencida la
prenotación, los herederos pueden enajenar libremente los bienes.
\end{formula}

\begin{important}
Son cinco años desde el \textbf{día presuntivo del fallecimiento}, no
desde el día en que la persona desapareció, y no desde la fecha de la
sentencia. Por eso el día presuntivo es tan relevante: es el que pone en
marcha el plazo de prenotación.
\end{important}

\section{Efectos de la reaparición (Art. 91 CCCN)}

Si el ausente reaparece \textbf{antes} de que venza la prenotación: se le
devuelven todos los bienes, porque el inventario asegura que se sabe
exactamente con qué se empezó, y la restricción de enajenación junto con
el requisito de autorización judicial protegieron el patrimonio.

Si el ausente reaparece \textbf{después} de vencida la prenotación
(cuando ya los bienes podían enajenarse libremente):

\begin{itemize}
    \item Se le entregan los bienes que aún existen.
    \item Se le entregan los bienes adquiridos con el valor de los que faltaron (subrogación real: por ejemplo, si se vendió un departamento y se compraron cuatro cocheras, se entregan las cuatro cocheras).
    \item Se le entrega el precio pendiente de cobro si se vendió un bien con saldo impago.
    \item Se le entregan los frutos no consumidos.
\end{itemize}

\begin{example}{Caso real: Oberá, Misiones (2011)}
Una persona se metió al río Uruguay y desapareció; la buscaron por tierra
y aire con helicópteros. Tiempo después, publicó en una red social una
foto desde Curitiba, en la playa. Se armó un escándalo: resultó que había
decidido hacerse ``desaparecer'' voluntariamente y fue descubierto con
una nueva familia. Si esa persona hubiera tenido bienes y se hubiera
iniciado la declaración, podría haber reaparecido y reclamado la
devolución de lo que le correspondía.
\end{example}

\begin{example}{Caso real: Inglaterra — reaparición tras cinco años}
Se conoció un caso en Inglaterra (noticias de diciembre de 2007) de una
persona que había desaparecido y reapareció cinco años después. Si
durante ese tiempo se hubiera declarado su fallecimiento presunto en
Argentina, al reaparecer tendría derecho a que se le restituyeran los
bienes —o su equivalente— conforme las reglas del Art. 91 CCCN.
\end{example}

% ======================================================================
% CUADRO CORNELL
% ======================================================================
\chapter{Cuadro de repaso Cornell}

El siguiente cuadro reúne, en formato Cornell, las preguntas y conceptos
centrales desarrollados a lo largo de los capítulos anteriores, con el
propósito de facilitar el repaso activo del material.

\begin{longtable}{
    >{\raggedright\arraybackslash}p{0.30\textwidth}
    >{\raggedright\arraybackslash}p{0.58\textwidth}
}
\toprule
\textbf{Preguntas / palabras clave} & \textbf{Notas / respuestas} \\
\midrule
\endfirsthead

\toprule
\textbf{Preguntas / palabras clave} & \textbf{Notas / respuestas} \\
\midrule
\endhead

\midrule
\multicolumn{2}{r}{\textit{Continúa en la página siguiente}} \\
\endfoot

\bottomrule
\endlastfoot

\textbf{¿Cuál es la diferencia entre ausencia y presunción de fallecimiento?} &
La ausencia (Arts. 79-84) es una medida conservatoria de bienes: nombra
un curador sin declarar fallecida a la persona. La presunción de
fallecimiento (Arts. 85-92) declara judicialmente el fallecimiento,
inscribe la sentencia en el Registro Civil y abre la sucesión. La
presunción supone ausencia, pero no es lo mismo. \\

\textbf{¿Quién puede pedir la declaración de ausencia?} &
El Ministerio Público y las personas que tengan interés legítimo en los
bienes (acreedores, sucesores). El juez competente es el del domicilio
del ausente; si no tiene domicilio en el país o es desconocido, el juez
del lugar donde están los bienes. \\

\textbf{¿Cuáles son los tres casos de presunción de fallecimiento?} &
1) Ordinario: ausencia sin noticias por 3 años (Art. 85). 2)
Extraordinario genérico: persona involucrada en catástrofe o actividad de
riesgo, plazo de 2 años (Art. 86, 1.\textsuperscript{er} párr.). 3)
Extraordinario específico: buque naufragado o aeronave perdida, plazo de
6 meses (Art. 86, 2.\textsuperscript{do} párr.). \\

\textbf{¿Cómo se fija el día presuntivo de fallecimiento?} &
En el ordinario: mitad del plazo (1 año y medio después de la última
noticia). En el extraordinario genérico: el día en que ocurrió el suceso
(o término medio si se extendió). En el extraordinario específico: el
último día en que se tuvo noticia del buque o aeronave. \\

\textbf{¿Cuáles son las tres garantías que protegen al ausente por si
reaparece?} &
1) Inventario detallado de todos los bienes. 2) Prenotación registral en
los asientos de los bienes. 3) Prohibición de enajenar o gravar sin
autorización judicial. La prenotación dura 5 años desde el día presuntivo
o hasta que la persona hubiera cumplido 80 años. \\

\textbf{¿Qué pasa si el ausente reaparece después de vencida la
prenotación?} &
Se le restituyen los bienes existentes, los adquiridos con el valor de
los que faltaron (subrogación real), el precio pendiente de cobro y los
frutos no consumidos. No puede recuperar bienes ya enajenados libremente
por los herederos una vez caduca la prenotación. \\

\textbf{¿Es necesario haber tramitado la ausencia antes de iniciar la
presunción de fallecimiento?} &
No. Según el Código, la declaración de ausencia (Art. 79) no es
presupuesto necesario ni suficiente para la presunción de fallecimiento.
Pueden tramitarse en forma sucesiva, pero en el juicio de presunción hay
que probar nuevamente las diligencias para dar con el ausente. \\

\textbf{¿Qué regula el CCCN sobre el procedimiento?} &
El CCCN incorpora normas procesales comunes para todo el país (citación
por edictos una vez por mes durante 6 meses, defensor oficial del
ausente, curador de bienes), para evitar diferencias entre jurisdicciones
en situaciones de gran sensibilidad jurídica y personal. \\

\midrule

\multicolumn{2}{p{\textwidth}}{
\textbf{Resumen:} El derecho civil argentino responde a la incertidumbre
generada por la desaparición de una persona mediante dos institutos
diferenciados que conviven en el CCCN (Arts. 79-92). La \textbf{ausencia}
(Arts. 79-84) es una herramienta de protección patrimonial preventiva:
permite nombrar un curador de bienes y adoptar medidas conservatorias sin
declarar fallecida a la persona. La \textbf{presunción de fallecimiento}
(Arts. 85-92) va más lejos: declara judicialmente el fallecimiento,
inscribe la sentencia en el Registro Civil y habilita la apertura de la
sucesión. Se estructura en tres casos según el grado de probabilidad de
muerte: el ordinario (3 años), el extraordinario genérico —catástrofes o
actividades de riesgo— (2 años) y el extraordinario específico —buque o
aeronave— (6 meses). La sentencia fija el día presuntivo, que determina
quiénes son los herederos y desde cuándo corre la prenotación registral
de cinco años. Durante ese período, los bienes no pueden enajenarse sin
autorización judicial, garantizando la restitución al ausente en caso de
reaparición. El sistema, con sus antecedentes en Vélez Sársfield y la Ley
14.394, busca equilibrar la certeza jurídica que exigen las relaciones
patrimoniales y sucesorias con la protección de los derechos del
ausente, a quien siempre se le reserva la posibilidad de recuperar su
patrimonio si regresa.
} \\

\bottomrule
\end{tabularx}
\end{table}

\end{document}
